\chapter{Conclusion and Future Scope}

This chapter draws together the main findings from our work, acknowledges the current limitations, and outlines the directions we see for extending the system.

\section{Conclusion}

We set out to build a system that could take over the repetitive parts of GitHub issue triage---classifying issues, tracing root causes to specific code locations, and proposing fixes---while remaining secure and scalable. Ansieyes achieves this through a combination of \gls{llm}-powered analysis, multi-layered security, and a two-pass architecture for large codebases.

The core analyser, powered by Gemini 2.0 Flash, classifies issues with about 91\% accuracy and produces root cause analyses that correctly point to the relevant source files in 85\% of cases. The quality validation and retry mechanism ensures that virtually every issue receives at least one actionable solution proposal. Duplicate detection works at two speeds: the Gemini-based approach catches semantic duplicates with an F1 of 0.87, while the cosine similarity fallback runs near-instantly for rapid pre-screening.

On the security front, the three-layer detection pipeline catches 96\% of prompt injection attempts across seven attack categories. The conservative aggregation strategy---taking the maximum risk from any detector---means the system errs on the side of caution, which we consider the right trade-off for an automated pipeline processing untrusted input.

The Librarian--Surgeon architecture proved essential for practical deployment. On repositories with 5,000 to 100,000 lines of code, it reduced token consumption by 57--70\%, and for larger codebases it enabled analysis that the single-pass approach simply could not handle.

\section{Future Scope}

Several extensions would strengthen the system:

\begin{enumerate}
\item \textbf{Multi-model support.} Adding backends for GPT-4, Claude, and open-source models (Llama, Mistral) would let teams choose based on cost, privacy, or performance needs.

\item \textbf{Multi-lingual issues.} Many open-source projects have contributors who file issues in languages other than English. Gemini already supports multiple languages, so adapting the prompts and evaluation pipeline would be the main work.

\item \textbf{Feedback-driven improvement.} When a developer edits or overrides an analysis, that feedback could be collected and used to refine prompt templates over time, creating a lightweight learning loop without full model fine-tuning.

\item \textbf{Platform integration.} Extending support to GitLab, Jira, and Azure DevOps would broaden the system's applicability beyond GitHub.

\item \textbf{Automated fix generation.} The current system proposes solutions in text form. A natural next step is generating actual pull requests with the proposed code changes, subject to human review before merging.

\item \textbf{Adversarial robustness.} The prompt injection detector could benefit from periodic red-teaming exercises and the addition of adversarial training data to keep pace with evolving attack techniques.
\end{enumerate}

\section{Learning Outcomes of the Project}
\begin{itemize}
\item Gained hands-on experience with \gls{llm} integration, including prompt design, structured output parsing, and handling model limitations like hallucination.
\item Built a multi-layered security system from scratch, understanding how \gls{ml}-based detection, pattern matching, and heuristics complement each other.
\item Developed proficiency with Pydantic for runtime type checking, pytest for structured testing, and GitHub Actions for \gls{ci}/\gls{cd} automation.
\item Learned to design for scalability constraints, particularly the token limits of \gls{llm} context windows, through the Librarian--Surgeon two-pass approach.
\item Worked with \gls{nlp} techniques (\gls{tfidf} vectorisation, cosine similarity) and understood their trade-offs against \gls{llm}-based alternatives.
\item Gained experience packaging Python tools for distribution, including \gls{cli} design and configuration management.
\item Understood the practical realities of deploying \gls{ai} systems: \gls{api} cost management, error handling, retry strategies, and the importance of validating model outputs before acting on them.
\end{itemize}
